% Propietario conservado: /Users/ruben/Documents/New project/output/REVISION_ARGUMENTAL_APERTURAS_CIERRES_20260915/VIII/spanish_source/ampliacion/schur_y_coercividad.tex
\clearpage
\section{Réduction de Schur et coercivité de la forme de Weil}
\label{viii:sec:schur-coercividad}
% Fuente íntegra: ampliacion/expediente/AGENTE_SCHUR_DETALLES.md

% PROCEDENCIA 11 de septiembre de 2026. Desarrollo focal para integración por el editor único.
% PROCEDENCIA No modifica los manuscritos I–VIII ni sus PDF.

\subsection{Base, domaine et résultat obtenu}
\label{viii:schur-add-sub-6}

On utilise les colonnes complètes déjà construites à partir de l'arithmétique native
d'APP, de l'état orienté TRIT, des opérations et de la mémoire TPK, et de la
structure conjointe du continu. Le produit natif publie des irréductibles et des
puissances ; leurs longueurs et leurs poids alimentent ensuite les lecteurs premier,
gamma et polaire. La coordonnée interne de \(\pi\) et le lecteur rationnel de la partie
finie d'Euler--Mascheroni interviennent après cette publication. Aucun
zéro, aucune valeur recherchée de la fonctionnelle ni aucune constante métrologique ne sélectionne ces
opérations.

On maintient les colonnes complètes et les domaines construits dans le
chapitre du double cercle et du transport spectral, avec les estimations
des supports et des détails. Les colonnes et le connecteur sont également explicités dans
la section~\ref{viii:sec:lectores-gamma}. La construction conjointe
du continu demeure un antécédent commun : elle n'est pas remplacée par cette
décomposition fonctionnelle, ni ses cinq constructions par cinq modules.

La nouvelle opération réunie ici est l'élimination exacte de l'espace infini des détails au moyen de sa forme coercive. Elle produit une matrice finie avec toutes les corrections de couplage, et pas seulement une matrice de moyennes. En outre, un cas concret est résolu :

Pour toute fonction de l'adhérence de
\(C_c^\infty((-1/18,1/18))\) dans la norme conjointe des lecteurs,
\begin{equation}
\mathscr W(f,f)\geq\frac12\|f\|_2^2.
\label{viii:schur-add-display-1}
\end{equation}
La même affirmation vaut dans tout intervalle translaté de longueur \(1/9\), et dans tous ses raffinements intérieurs nonadiques.

La fenêtre complète a pour longueur \(1/9>0.085626\ldots\), l'extrémité
du critère basal précédent. Elle ne contient aucune translation première active,
car \(1/9<\log2\). Cette limite de l'exemple reste explicite.

\textbf{Provenance.} Architecture de l'auteur et lecteurs : préexistants. Coercivité
des détails : résultat récupéré du théorème de seuil pour la moyenne nulle, en conservant une estimation
intermédiaire plus forte que sa borne logarithmique abrégée. Assemblage de Schur,
certification résiduelle et preuve de l'intervalle \(1/9\) : formalisation et
certificat ajoutés dans cette section ; aucune nouveauté historique globale n'est revendiquée.

\subsection{Partition, moyennes cellulaires et détails}
\label{viii:schur-add-sub-46}

Fixons \(R>0\), \(I_R=(-R/2,R/2)\) et \(H_R=L^2(I_R)\), avec prolongement par zéro. Soit \(V_R\) l'adhérence de \(\mathcal D_R=C_c^\infty(I_R)\) pour la norme conjointe des lecteurs complets. De manière équivalente, sa norme peut être prise comme
\begin{equation}
\begin{aligned}
\|f\|_{V_R}^2&=\|f\|_2^2+Q_\Gamma(f),\\
Q_\Gamma(f)&=\int_0^\infty\mu(a)\|(I-T_a)f\|_2^2\,da,\\
\mu(a)&=\frac{e^{-a/2}}{1-e^{-2a}}.
\end{aligned}
\label{viii:schur-add-display-2}
\end{equation}
Les canaux premiers et polaires sont bornés dans chaque fenêtre fixée.
La forme gamma est fermée : c'est la forme d'un multiplicateur positif
en Fourier, restreinte au sous-espace des fonctions à support dans la
fenêtre ; de manière équivalente, cela se vérifie par convergence dans l'intégrale
directe de différences. La forme complète est une perturbation bornée
de celle-ci et définit un opérateur autoadjoint semi-borné \(L_R\).

Partitionnons \(I_R\) en intervalles contigus \(I_j=(\ell_j,r_j)\), de longueurs \(h_j\leq h\), et posons
\begin{equation}
u_j=h_j^{-1/2}\mathbf1_{I_j},\qquad
Ua=\sum_{j=1}^M a_j u_j,\qquad P=UU^*,\quad Q=I-P.
\label{viii:schur-add-display-3}
\end{equation}
Dans la dernière somme, \(a=(a_1,\ldots,a_M)\) est le vecteur des coefficients ;
les extrémités \(\ell_j,r_j\) ne sont utilisées que pour définir \(I_j\).

Les indicatrices appartiennent à \(V_R\) : leur énergie gamma est contrôlée par \(\|(I-T_a)u_j\|_2^2=2\min(|a|/h_j,1)\). Ainsi, \(P\) agit continûment dans \(V_R\), et tout \(f\in V_R\) possède une décomposition unique
\begin{equation}
f=Ua+d,\qquad a_j=\langle f,u_j\rangle,\qquad
d\in V_D:=V_R\cap QH_R.
\label{viii:schur-add-display-4}
\end{equation}
Ici, \(d\) a une intégrale nulle dans chaque cellule. L'espace \(V_D\) est
fermé dans la norme de forme et dense dans \(QH_R\).

L'estimation des détails de moyenne nulle donne une borne \(\mathscr W_R(d,d)\geq\eta\|d\|_2^2\) si l'on choisit une partition suffisamment fine. On peut employer son seuil logarithmique publié, ou conserver l'estimation intermédiaire exacte
\begin{equation}
\begin{aligned}
\eta_N(h,R)={}&4\sum_{n=0}^N\frac1{4n+1}
-\frac{h^2(N+1)(2N+1)}{\pi^2}\\
&+c_\Gamma-2\sum_{k\log p\leq R}(\log p)p^{-k/2}-\beta_R,\\
\beta_R={}&2\sinh(R/2)-R.
\end{aligned}
\label{viii:schur-add-eq-S1}
\end{equation}
\begin{lemma}[Borne retenue pour les détails cellulaires]
\label{viii:schur-add-retencion-detalles}
Pour \(d\in V_D\) et tout entier \(N\geq0\),
\[
\mathscr W_R(d,d)\geq\eta_N(h,R)\|d\|_2^2.
\]
\end{lemma}
\begin{proof}
Dans chaque cellule \([\ell_j,r_j]\), on définit
\(F(x)=\int_{\ell_j}^x d(t)\,dt\). La moyenne nulle donne
\(F(\ell_j)=F(r_j)=0\). La réunion de ces primitives, prolongée
par zéro, appartient à \(H_0^1(I_R)\), satisfait \(F'=d\) et, par
Dirichlet dans chaque cellule,
\begin{equation}
\|F\|_2^2\leq\frac{h^2}{\pi^2}\|d\|_2^2.
\label{viii:schur-add-primitiva}
\end{equation}
Le développement positif
\(\mu(a)=\sum_{n\geq0}e^{-\lambda_n a}\), \(\lambda_n=2n+1/2\),
permet d'écrire \(Q_\Gamma=\sum E_{\lambda_n}\), où
\begin{equation}
E_\lambda(d)=\int_0^\infty e^{-\lambda a}
 \|(I-T_a)d\|_2^2\,da
 =\frac2\lambda\|d\|_2^2-\langle d,C_\lambda d\rangle\geq0,
\label{viii:schur-add-energia-lambda}
\end{equation}
et \(C_\lambda\) est la convolution par \(e^{-\lambda|x|}\).
Pour la transformée unitaire
\(\mathcal F_u=(2\pi)^{-1/2}\widehat{\phantom f}\), le multiplicateur
de \(C_\lambda\) est \(2\lambda/(\lambda^2+\xi^2)\). Par Plancherel,
\begin{equation}
\begin{aligned}
\langle d,C_\lambda d\rangle
&=\int_{\mathbb R}\frac{2\lambda\xi^2}{\lambda^2+\xi^2}
 |\mathcal F_uF(\xi)|^2\,d\xi\\
&\leq2\lambda\|F\|_2^2
 \leq\frac{2\lambda h^2}{\pi^2}\|d\|_2^2 .
\end{aligned}
\label{viii:schur-add-cota-convolucion}
\end{equation}
Retenir \(N+1\) termes et écarter uniquement la queue non négative donne
\[
Q_\Gamma(d)\geq
\left[4\sum_{n=0}^N\frac1{4n+1}
-\frac{h^2(N+1)(2N+1)}{\pi^2}\right]\|d\|_2^2.
\]
Le terme scalaire apporte \(c_\Gamma\|d\|_2^2\). Chaque translation première apporte au moins \(-2w_{p,k}\|d\|_2^2\), par Cauchy--Schwarz et l'unitarité de \(T_{k\log p}\). Seules interviennent les longueurs inférieures ou égales à \(R\). Enfin, la contribution polaire est
\[
2\left|\int_{I_R}d(x)\cosh(x/2)\,dx\right|^2
-2\left|\int_{I_R}d(x)\sinh(x/2)\,dx\right|^2
\geq-\beta_R\|d\|_2^2,
\]
car \(2\int_{I_R}\sinh^2(x/2)\,dx=2\sinh(R/2)-R\).
La somme prouve la borne.
\end{proof}
Il n'est pas nécessaire de remplacer la somme harmonique par son intégrale inférieure.
Pour le certificat de la fenêtre \(1/9\), on vérifie en outre
\((2N+1/2)h<\pi\), de sorte que chaque borne gamma retenue est non négative.

Pour tout \(R\) fixé, il existe une partition qui rend positive cette
borne. En effet, si \(0<h\leq1\) et \(N=\lfloor1/h\rfloor\), alors
\[
4\sum_{n=0}^N\frac1{4n+1}\geq\log(4N+5)\geq\log(1+4/h),
\qquad
\frac{h^2(N+1)(2N+1)}{\pi^2}<\frac23.
\]
Le premier terme diverge lorsque \(h\downarrow0\) ; les termes
premier, scalaire et polaire restent fixés par \(R\). Cela conserve
la distinction entre coercivité des détails et signe de la
forme complète.

\subsection{Le complément de Schur exact est de taille finie}
\label{viii:schur-add-sub-97}

Soit \(D\) l'opérateur associé à la restriction de la forme à \(V_D\).
Si \(\eta>0\), alors \(D\geq\eta I\) sur \(QH_R\),
\(D^{-1}\) existe et \(\|D^{-1}\|\leq\eta^{-1}\).
On ne présume pas que \(L_R\) conserve l'espace des détails.

Les indicatrices appartiennent également au domaine de l'opérateur \(L_R\). L'expression de leur image gamma est donnée dans la section suivante : elle ne possède que des singularités logarithmiques, qui appartiennent à \(L^2(I_R)\). Cela permet de définir, sans étendre indûment une fonctionnelle de forme,
\begin{equation}
A=U^*L_RU,\qquad B=QL_RU:\mathbb C^M\longrightarrow QH_R.
\label{viii:schur-add-display-6}
\end{equation}
Le produit scalaire est linéaire dans la première variable, comme dans le
manuscrit. Les matrices sont écrites de sorte que
\(\mathscr W_R(Ua,Ua)=a^*Aa\). L'identité par blocs est
\begin{equation}
\mathscr W_R(Ua+d,Ua+d)
=a^*Aa+2\operatorname{Re}\langle Ba,d\rangle
+\mathfrak d[d,d],
\label{viii:schur-add-display-7}
\end{equation}
où \(\mathfrak d\) est la forme fermée de \(D\).

\begin{theorem}[Élimination exacte de l'espace des détails]
\label{viii:schur-add-schur-exacto}
La matrice
\begin{equation}
S=A-B^*D^{-1}B
\label{viii:schur-add-eq-S2}
\end{equation}
satisfait l'identité exacte
\begin{equation}
\boxed{\mathscr W_R(Ua+d,Ua+d)
=a^*Sa+
\mathfrak d[d+D^{-1}Ba,d+D^{-1}Ba].}
\label{viii:schur-add-eq-S3}
\end{equation}
En particulier, \(\mathscr W_R\geq0\) sur \(V_R\) si et seulement si \(S\succeq0\). Son indice négatif coïncide avec celui de \(S\).
\end{theorem}
\begin{proof} Pour \(z=D^{-1}Ba\in\operatorname{Dom}D\),
\(\mathfrak d[z,d]=\langle Ba,d\rangle\). Développer le second terme
de \eqref{viii:schur-add-eq-S3} et utiliser \(\mathfrak d[z,z]=a^*B^*D^{-1}Ba\) donne l'égalité.
La suffisance est immédiate. Pour la nécessité, on choisit
\(d=-D^{-1}Ba\). Le changement triangulaire \((a,d)\mapsto(a,d+D^{-1}Ba)\)
est inversible dans \(\mathbb C^M\oplus V_D\). La seconde forme résultante
est définie positive ; les directions négatives sont exactement celles de
la matrice finie.
\end{proof}

Ce théorème transforme le résultat d'indice négatif fini en une matrice exacte. La matrice contient encore \(D^{-1}\) : la finitude de sa taille ne
signifie pas que ses entrées soient déjà calculées. La section suivante
donne une manière de les encadrer avec des résidus contrôlés.

\subsection{Images explicites et bornes certifiables des corrections}
\label{viii:schur-add-sub-148}

Posons \(M(x)=\int_x^\infty\mu(a)\,da\), pour \(x>0\). Pour une indicatrice \(u=h^{-1/2}\mathbf1_{(a,b)}\), son image gamma, restreinte à \(I_R\), est presque partout
\begin{equation}
(L_\Gamma u)(x)=h^{-1/2}
\begin{cases}
M(x-a)+M(b-x),&a<x<b,\\
-M(a-x)+M(b-x),&x<a,\\
-M(x-b)+M(x-a),&x>b.
\end{cases}
\label{viii:schur-add-eq-S4}
\end{equation}
L'égalité s'obtient en intégrant les deux translations : à l'intérieur de la
cellule, il reste une queue pour chaque extrémité, et à l'extérieur, il reste l'intégrale de l'
unique intervalle de translations qui touche la cellule. On peut la prouver d'abord
avec le régulateur \(a_{\rm traslado}\geq\epsilon\) ; la convergence dans
\(L^2(I_R)\) de l'expression logarithmique justifie sa limite faible de
forme. Les autres parties sont
\begin{equation}
L_{\rm primo}u=-\sum_{k\log p\leq R}w_{p,k}
M_R(T_{k\log p}+T_{-k\log p})u,
\label{viii:schur-add-display-11}
\end{equation}
\begin{equation}
\begin{aligned}
L_{\rm polar}u(x)&=2\int_a^b h^{-1/2}\cosh((x-y)/2)\,dy,\\
L_Ru&=L_\Gamma u+c_\Gamma u+L_{\rm primo}u+L_{\rm polar}u.
\end{aligned}
\label{viii:schur-add-eq-S5}
\end{equation}
Toutes les discontinuités et singularités de \eqref{viii:schur-add-eq-S4}--\eqref{viii:schur-add-eq-S5} se trouvent dans une liste finie d'extrémités et d'extrémités translatées par des horloges publiées. Ainsi, \(A_{ij}\) et
\begin{equation}
G_B=B^*B,\qquad (G_B)_{ij}=\langle Be_j,Be_i\rangle
\label{viii:schur-add-display-13}
\end{equation}
se calculent par des intégrales unidimensionnelles explicites. Dans cette dernière
formule, \(e_j\) désigne le vecteur de coordonnées standard de
\(\mathbb C^M\), et non l'indicatrice dans \(H_R\).
La borne initiale est
\begin{equation}
A-\eta^{-1}G_B\preceq S\preceq A.
\label{viii:schur-add-eq-S6}
\end{equation}
C'est une inégalité opératorielle dont les termes croisés sont conservés, et non une substitution des entrées diagonales de \(G_B\) à l'opérateur lui-même.

\begin{theorem}[Certification du complément par le résiduel]
\label{viii:schur-add-residual}
Pour tout opérateur fini d'approximation \(Z:\mathbb C^M\to\operatorname{Dom}D\), définissons
\begin{equation}
E=B-DZ,\qquad C_Z=Z^*B+B^*Z-Z^*DZ.
\label{viii:schur-add-display-15}
\end{equation}
Alors
\begin{equation}
B^*D^{-1}B=C_Z+E^*D^{-1}E,
\label{viii:schur-add-display-16}
\end{equation}
et donc
\begin{equation}
\boxed{A-C_Z-\eta^{-1}E^*E\preceq S\preceq A-C_Z.}
\label{viii:schur-add-eq-S7}
\end{equation}
\end{theorem}
\begin{proof}
Substituer \(B=DZ+E\) et développer le produit. Tous les
termes sont définis parce que les colonnes de \(Z\) appartiennent à
\(\operatorname{Dom}D\). La borne utilise \(0<D^{-1}\leq\eta^{-1}I\).
\end{proof}

Un choix effectif de \(Z\) consiste à résoudre le système de Galerkin sur des détails en escalier d'un raffinement nonadique. Ses colonnes sont des combinaisons finies d'indicatrices et ont les images \eqref{viii:schur-add-eq-S4}--\eqref{viii:schur-add-eq-S5}, de sorte que \(E\) et son Gram sont explicitement intégrables. Le système de Galerkin n'est pas déclaré solution exacte : le terme \(\eta^{-1}E^*E\) conserve son erreur sur le complément infini.

Pour certifier ces intégrales sans perdre les extrémités singulières,
on sépare
\begin{equation}
M(x)=-\tfrac12\log x+\log2+\tfrac\pi4+r(x),
\qquad r(0)=0,\quad r'(x)=\tfrac1{2x}-\mu(x).
\label{viii:schur-add-eq-S8}
\end{equation}
Dans \(0<x\leq R\),
\begin{equation}
-\tfrac14 e^{R/2}\leq r'(x)\leq\tfrac14,
\qquad |r'(x)|\leq\tfrac14e^{R/2}.
\label{viii:schur-add-display-19}
\end{equation}
En effet, \(\mu(x)=e^{x/2}/(2\sinh x)\) et
\(x\leq\sinh x\leq xe^x\) donnent
\(e^{-x/2}/(2x)\leq\mu(x)\leq e^{x/2}/(2x)\).
Les inégalités \(1-e^{-y}\leq y\) et \(e^y-1\leq ye^y\)
prouvent les bornes. On n'affirme pas que \(r'\) ait un signe constant
dans une fenêtre arbitrairement grande.
On intègre analytiquement \(\log x\) et \(\log^2x\) dans les petits
tronçons d'extrémité ; le reste \(r\) est borné par sa constante de Lipschitz.
Hors des extrémités, on utilise des quadratures par intervalles subdivisées et
la série exponentielle de \(M\). Par exemple,
\begin{equation}
\int_0^\delta\log^2x\,dx
=\delta[(\log\delta)^2-2\log\delta+2].
\label{viii:schur-add-display-20}
\end{equation}
Les produits de deux singularités en des points distincts sont traités
en séparant leurs voisinages ; pour une extrémité commune, on utilise la même
intégrale quadratique. Les intervalles résultants sont insérés dans \eqref{viii:schur-add-eq-S7} et
la positivité est vérifiée par une décomposition \(LDL^*\) avec
des pivots intervalles strictement positifs ou par une borne de perturbation
d'une matrice rationnelle. Un certificat fini est ainsi spécifié
avec un contrôle du complément infini ; on n'utilise pas de maillage sans borne de queue.

\subsection{Coercivité dans une fenêtre de longueur 1/9}
\label{viii:schur-add-sub-248}

Prenons \(R=h=1/9\), une seule cellule et \(u=R^{-1/2}\mathbf1_{I_R}\). C'est un intervalle complet, et non neuf intervalles séparés par des lacunes. La section des détails est tout \(V_R\cap\{u\}^{\perp}\).

\subsubsection{Marge du bloc infini}
\label{viii:schur-add-sub-254}

Dans \eqref{viii:schur-add-eq-S1}, on prend \(N=13\). Il n'y a pas de termes premiers puisque \(R<\log2\). On obtient
\begin{equation}
\eta=4\sum_{n=0}^{13}\frac1{4n+1}
-\frac{(1/9)^2\,14\cdot27}{\pi^2}+c_\Gamma-\beta_{1/9}>1.
\label{viii:schur-add-eq-S9}
\end{equation}
Le certificat par intervalles donne
\(1.00363976543713133046<\eta<1.00363976543713133048\).
La condition \((53/2)/9<\pi\) garantit en outre la positivité de
chaque borne gamma retenue. La somme est conservée ; la remplacer par la
borne \(\log(1+4/h)-2/3\) ferait perdre ici la marge nécessaire.

\subsubsection{Bloc de moyenne}
\label{viii:schur-add-sub-269}

À l'intérieur de la fenêtre, en écrivant \(x=R(t-1/2)\),
\begin{equation}
L_Ru(x)=R^{-1/2}g(t),\quad
g(t)=c_\Gamma+M(Rt)+M(R(1-t))
+8\sinh(R/4)\cosh(R(t-1/2)/2).
\label{viii:schur-add-display-22}
\end{equation}
Ainsi, \(A=\int_0^1g(t)\,dt\) et \(\|B\|^2\) est la variance
de \(g\) sur l'intervalle unité. Tonelli appliqué à \(M\) donne
\begin{equation}
A=c_\Gamma+\frac2R\sum_{n=0}^\infty
\frac{1-e^{-(2n+1/2)R}}{(2n+1/2)^2}
+\frac{32}R\sinh^2(R/4).
\label{viii:schur-add-eq-S10}
\end{equation}
Si l'on retient les \(L\) premiers termes, la queue positive vaut au plus \(\lambda_L^{-2}+(2\lambda_L)^{-1}\), avec \(\lambda_L=2L+1/2\). On sépare le premier terme de la queue et on majore le reste par l'intégrale d'une fonction décroissante. Pour \(L=1024\), le certificat produit
\begin{equation}
0.9723307136651774<A<0.9767284617331572,
\qquad A>97/100.
\label{viii:schur-add-eq-S11}
\end{equation}
Une approximation plus précise n'est pas nécessaire pour la preuve.

\subsubsection{Couplage entre moyennes et détails}
\label{viii:schur-add-sub-295}

Dans \(0<x\leq1/9\), l'estimation de \eqref{viii:schur-add-eq-S8} est améliorée en
\begin{equation}
-9/10\leq r'(x)\leq0.
\label{viii:schur-add-display-25}
\end{equation}
En effet, pour \(0<x\leq1/9\),
\begin{equation}
\frac{\sinh x}{x}
=\sum_{k\geq0}\frac{x^{2k}}{(2k+1)!}
\leq\frac1{1-x^2}\leq1+\frac x2\leq e^{x/2}.
\label{viii:schur-add-display-26}
\end{equation}
L'inégalité intermédiaire équivaut à \(1-2x-x^2\geq0\),
qui vaut sur cet intervalle. Ainsi, \(\mu(x)\geq1/(2x)\).
D'autre part,
\begin{equation}
\mu(x)-\frac1{2x}\leq\frac{e^{x/2}-1}{2x}
\leq\frac14 e^{x/2}
\leq\frac1{4(1-1/18)}=\frac9{34}<\frac9{10}.
\label{viii:schur-add-display-27}
\end{equation}
On conserve la borne \(9/10\), moins précise, utilisée par le certificat.
La fonction \(q(t)=r(Rt)+r(R(1-t))\) est symétrique par rapport à
\(1/2\) et sa constante de Lipschitz est au plus \(9R/10\).
Son oscillation est au plus \(9R/20=1/20\) ; son écart-type
est donc au plus \(1/40\). Cette dernière affirmation découle de
\(\operatorname{Var}(X)\leq(\sup X-\inf X)^2/4\) : soustraire le
milieu de l'étendue et utiliser le fait que la moyenne minimise l'erreur quadratique.

La partie singulière \(v(t)=-\tfrac12\log(t(1-t))\) a pour moyenne 1 et pour variance \(1-\pi^2/12\). Pour le vérifier, on intègre \(\log t\), \(\log^2t\) et la série positive de \(\log t\log(1-t)\), qui donne \(\int_0^1\log t\log(1-t)\,dt=2-\sum_{n\geq1}n^{-2}\). L'identité \(\sum n^{-2}=\pi^2/6\) peut s'obtenir par Parseval appliqué à la série de sinus de \(x\) sur \((-\pi,\pi)\) : ses coefficients sont \(2(-1)^{n+1}/n\) et \((1/\pi)\int_{-\pi}^{\pi}x^2dx=2\pi^2/3\). C'est une reconnaissance fonctionnelle ultérieure de la même \(\pi\), et non une sélection de son générateur.

La partie polaire a une oscillation
\(8\sinh(R/4)[\cosh(R/4)-1]\) et un écart-type au plus égal à
sa moitié. L'inégalité triangulaire dans l'espace des fonctions
centrées donne, pour le couplage complet,
\begin{equation}
\boxed{\|B\|^2\leq
\left[\sqrt{1-\pi^2/12}+\frac1{40}
+4\sinh(R/4)(\cosh(R/4)-1)\right]^2<\frac15.}
\label{viii:schur-add-eq-S12}
\end{equation}
Le terme scalaire et toutes les constantes indépendantes de \(t\)
s'éliminent par centrage, et non par une prétendue disparition des
termes croisés. L'évaluation par intervalles du membre majorant est
\(0.19926357337268538427\ldots<1/5\).

\subsubsection{Schur positif et coercivité de la forme entière}
\label{viii:schur-add-sub-348}

\begin{theorem}[Coercivité dans la fenêtre nonadique complète]
\label{viii:schur-add:coercividad-ventana}
Dans \(I_{1/9}=(-1/18,1/18)\), toute fonction du domaine conjoint
\(V_{1/9}\) satisfait
\[
\mathscr W_{1/9}(f,f)\geq\tfrac12\|f\|_2^2.
\]
La borne vaut sur tout intervalle translaté de même longueur et sur tous
ses raffinements intérieurs.
\end{theorem}
\begin{proof}
De \eqref{viii:schur-add-eq-S6}, \eqref{viii:schur-add-eq-S9}, \eqref{viii:schur-add-eq-S11} et \eqref{viii:schur-add-eq-S12},
\begin{equation}
S=A-B^*D^{-1}B>\frac{97}{100}-\frac15=\frac{77}{100}>0.
\label{viii:schur-add-display-29}
\end{equation}
De plus, pour \(f=au+d\),
\begin{equation}
\mathscr W_R(f,f)-\tfrac12\|f\|_2^2
\geq(A-\tfrac12)|a|^2
-2\|B\||a|\|d\|_2+(\eta-\tfrac12)\|d\|_2^2.
\label{viii:schur-add-display-30}
\end{equation}
La matrice scalaire du second membre a une diagonale positive et un déterminant strictement supérieur à \(\tfrac{47}{100}\cdot\tfrac12-\tfrac15=\tfrac7{200}>0\). On prouve ainsi \(\mathscr W_R(f,f)\geq\tfrac12\|f\|_2^2\) pour tout \(f\in V_R\). L'argument contrôle tous les détails infinis ; il ne les remplace pas par des indicatrices finies. La translation conserve les corrélations et multiplie les évaluations polaires par \(e^{t/2}\) et \(e^{-t/2}\), dont les produits croisés se compensent. Elle préserve donc également la forme complète et la borne.
\end{proof}

\subsection{Raffinement, double cercle et les deux moments}
\label{viii:schur-add-sub-367}

Chaque subdivision en neuf conserve exactement
\(u_{\rm padre}=\tfrac13\sum_{r=0}^8u_r\), et la recomposition des
fonctions intervient avant l'application de \(J_+\) et de \(J_-\). Par conséquent,
les gramiens de toutes les indicatrices nonadiques dans la fenêtre
satisfont \(G_n\geq\tfrac12 I\) et
\(V_n^*G_{n+1}V_n=G_n\). Les huit détails par parent et
tous leurs termes croisés sont conservés. Il s'agit d'une uniformité en profondeur sur cette fenêtre,
et non d'un résultat uniforme lorsque \(R\) augmente.

Pour une partition générale, la matrice de Schur exacte n'obéit pas simplement
à la même compression que le Gram brut. Si l'on raffine et exprime le
Schur fin en coordonnées ancienne moyenne/nouveau détail, le Schur
ancien s'obtient en éliminant également ce nouveau détail. C'est
l'associativité de la minimisation dans \eqref{viii:schur-add-eq-S3}. Affirmer
\(S_{\rm padre}=V^*S_{\rm hijos}V\) omettrait cette seconde correction.

Dans la fenêtre certifiée, le connecteur déjà construit \(\mathcal C_R=J_{-,R}L_R^{\rm cola}\), où \(L_R^{\rm cola}\) est la cotrace \eqref{viii:gamma-add-eq-11} appliquée à la composante gamma, satisfait \(\mathcal C_RJ_{+,R}=J_{-,R}\). Sa restriction à l'adhérence de l'image effective \(\mathcal M_{+,R}\) est contractante, puisque
\begin{equation}
\|J_{-,R}f\|^2\leq\|J_{+,R}f\|^2.
\label{viii:schur-add-display-31}
\end{equation}
La contraction se démontre maintenant par Schur ; elle n'est pas définie à
partir d'un Gram supposé positif. On prend ensuite son défaut
\(\Delta_R=(I-\mathcal C_R^*\mathcal C_R)^{1/2}\) dans cet espace
effectif et on construit
\begin{equation}
\Xi_Rf=J_{-,R}f\oplus\Delta_RJ_{+,R}f.
\label{viii:schur-add-display-32}
\end{equation}
Si \(P_-\) est la projection sur le premier terme de la somme,
\begin{equation}
\Xi_R^*\Xi_R=J_{+,R}^*J_{+,R},\qquad
\Xi_R^*P_-\Xi_R=J_{-,R}^*J_{-,R}.
\label{viii:schur-add-display-33}
\end{equation}
Ce sont des égalités de formes sur le domaine conjoint, et non des produits d'
opérateurs non bornés sur tout \(L^2\). Elles donnent les deux moments dans
le codomaine de dilatation. Une identification supplémentaire avec un
\(T_N\) concret du double cercle nécessite l'entrelacement
avec cet opérateur : on n'identifie pas une projection arbitraire à
\(T_N\) du seul fait de disposer des deux moments.

\subsection{Contrôles de reproduction et conditions d'extension}
\label{viii:schur-add-sub-411}

Le programme de vérification conservé dans l'annexe~\ref{viii:ap:procedencia-reproduccion} certifie les constantes par arithmétique d'intervalles et bornes explicites de queue. Il utilise le préfixe interne de \(\pi\) avec ses empreintes et le producteur rationnel déjà conservé de \(\gamma\) et \(\log2\). Il n'utilise pas \texttt{\detokenize{mpmath.pi}}, \texttt{\detokenize{mpmath.euler}}, de tables de zéros ni d'arguments flottants comme entrées.

% DOCUMENTACION ```text
% DOCUMENTACION python3 verificar_schur_ventana_nonadica.py --receipt CERTIFICADO_SCHUR_R_UN_NOVENO.json
% DOCUMENTACION python3 -O verificar_schur_ventana_nonadica.py --receipt CERTIFICADO_SCHUR_R_UN_NOVENO_OPTIMIZADO.json
% DOCUMENTACION ```

Les deux exécutions satisfont aux mêmes 21 contrôles par intervalles.
% ESTADO HISTORICO: PASS_SCHUR_FIXED_WINDOW_ONE_NINTH.
Les reçus conservent les extrémités rationnelles exactes, la précision, la version
de la bibliothèque, la provenance et les empreintes. Le contrôle adversarial
\(\begin{pmatrix}1&2\\2&1\end{pmatrix}\) possède deux blocs
diagonaux positifs et un Schur \(-3\) : il rejette l'utilisation de leurs seuls signes.

Falsificateurs substantiels : omettre le terme négatif polaire ; intervertir la moyenne
de cellule et l'échantillonnage ponctuel ; perdre une translation première dans une
fenêtre plus grande ; remplacer le Schur par \(A\) ; accepter Galerkin sans
\(E^*E/\eta\) ; ou confondre tous les raffinements intérieurs avec
toutes les fenêtres. Aucun ne fait partie de ce certificat.

L'extension concrète suivante est une fenêtre \(R>\log2\), avec
une partition nonadique dont le bloc de détails soit coercif. Il faut y
évaluer \eqref{viii:schur-add-eq-S7}, en incluant au moins l'horloge \(\log2\).
Le résultat présent n'affirme ni ce nouveau signe, ni la positivité
globale de Weil, ni RH. Il ne tire pas non plus de négation du corpus à partir
de cette portée locale : il fournit une élimination exacte, une procédure
certifiable et un premier cas clos qui étend le cas basal précédent.

% DOCUMENTACION ## Localizadores materiales
% DOCUMENTACION 
% DOCUMENTACION - `output/ARTICULO_VIII_PRIMOS_ESTRUCTURA_ESPECTRAL_20260910/manuscrito/03_DOBLE_CIRCULO_Y_TRANSPORTE_ESPECTRAL.md`: §§2–3, líneas 120–260; §12, 1000–1244; §13, 1245–1514; §14, 1515–1739; §15, 1740–1896.
% DOCUMENTACION - Mismo paquete, `fuente/pruebas/python/verificar_gram_nueve_ventanas.py`: publicación autenticada de π, forma intervalar y cola gamma.
% DOCUMENTACION - Mismo paquete, `fuente/pruebas/python/partes_finitas_exactas.py`: productores racionales de γ y logaritmos.
% DOCUMENTACION - Mismo paquete, `fuente/pruebas/python/verificar_nueve_intervalos_coercividad.py`: estimación completa previa en soportes desconectados.
% DOCUMENTACION - Mismo paquete, `gestion/recibo_gamma/RECIBO_CONSTANTES_GAMMA.json` y `RECIBO_GENEALOGIA_GAMMA.json`: procedencia heredada; esta sección no cambia su estado ni pretende estar ya incorporada en esos recibos.
% DOCUMENTACION - Serie activa: integral de 2.249 páginas y desarrollo estructural común. El gate global de primera aplicación sigue devolviendo el testigo de 2.084 páginas; conforme a continuidad, no desplaza la selección autoral de la serie ni este delta focal.
% DOCUMENTACION
